\SourceLeaf{275}%
selbstständige Geometrie, der er bei einer Gelegenheit nachgegangen sei und die er als antieuklidische Geometrie bezeichnet habe.

\textsc{Gauss} betrachtete nach seiner oft ausgesprochenen innersten Überzeugung die drei Dimensionen des Raums als eine besondere Eigenschaft der menschlichen Seele. Menschen, die das nicht einsehen konnten, nannte er einmal in humoristischer Laune Böotier. Wir können uns, sagte er, etwa in Wesen hineinversetzen, die sich nur zweier Dimensionen bewusst sind. Wesen, die höher als wir stehen, würden vielleicht auf ähnliche Weise auf uns herabblicken. Er habe, fuhr er scherzend fort, hier gewisse Probleme beiseitegelegt, die er später in einem höheren Zustand geometrisch behandeln wolle.\}

\Needspace{6\baselineskip}
\begin{center}\textsc{Bemerkung.}\end{center}

Die von \textsc{Sartorius von Waltershausen} erwähnte Schrift hat sich in \textsc{Gauss}’ Nachlass gefunden. Sie bezieht sich im Wesentlichen nur auf die Erklärung des Zahlbegriffs und der vier Grundrechenarten.

\hfill \textsc{Stäckel}.
